\chapter {Work Evaluation and Future Improvement}
\noindent
\rule{\textwidth}{0.4pt}

\begin{center}
\begin{minipage}{0.8\textwidth}
\textit{In Chapter 6, the study examines the key principles and practices that underpin continuous improvement, while also assessing the outcomes of the project within the context of an evolving technological landscape.}
\end{minipage}
\end{center}

\section{Work Evaluation}

Throughout the implementation process, the team organized the workflow in a clearly modular manner, ensuring effective task allocation among members and maintaining steady progress through short development cycles (sprints). An iterative approach was adopted to decompose complex requirements—such as the code runner, sandboxed execution, ranking system, AI chatbot, backend/API, and user interface—into manageable components. Each component was incrementally developed and reviewed at the end of every sprint to refine priorities and address issues. Design documentation and visual artefacts—including architectural diagrams, ERD, service flows, mockups, and progress charts—were thoroughly prepared, enabling instructors and stakeholders to effectively monitor progress and provide constructive feedback.

From a technical perspective, the team explored and implemented several system design techniques, including component decomposition, API design with formal documentation (Swagger/OpenAPI), and the development of a secure code execution environment (sandbox) with resource constraints and isolation mechanisms. The integration of an AI chatbot to enhance user experience provided valuable insights into incorporating NLP models within real-world application workflows. Mockups and prototypes facilitated early-stage usability testing, while architectural diagrams and flow representations ensured consistency between the design and implementation phases.

Continuous testing and evaluation (unit and integration testing, code runner execution flows, and API validation) were conducted in parallel with feature development to maintain system quality. The team also gained important experience in time management, team coordination, and resource allocation—particularly when addressing challenges related to execution environment security and performance optimization. In conclusion, the project successfully achieved its initial design objectives and demonstrated the system’s operational feasibility, while also establishing a solid foundation for further refinement, scalability, and real-world deployment.

\section{Future Improvement}

Due to time and resource constraints, the system has reached a functional and stable state, providing a solid foundation for further development. However, there remain several opportunities for improvement that can enhance overall system capability, user experience, and scalability in future iterations.

\textbf{Chatbot Enhancement:} The current agent supports coding interview practice by explaining algorithmic topics, recommending problems, listing solved or unsolved problems, and assisting users in debugging failed submissions. In the future, the system can be improved by integrating the agent with additional services such as contest management, notification systems, and user progress analytics. The agent could also provide more personalized learning paths based on submission history, solved problems, weak topics, and difficulty progression. Another potential improvement is to develop a separate moderator chatbot that assists moderators in managing problems, reviewing submissions, handling user reports, and generating system-wide notifications more efficiently.

\textbf{Code Runner Optimization and Scalability:} The current code execution system can be improved in terms of performance and scalability. Enhancements such as distributed execution, job queuing mechanisms, and more efficient resource allocation would allow the system to handle a larger number of concurrent submissions. This is particularly important as the user base grows and the demand for real-time feedback increases.

\textbf{Sandbox Security Enhancement:} Ensuring secure code execution remains a critical concern. Future improvements may include stricter isolation mechanisms, enhanced monitoring of resource usage, and protection against malicious code patterns. Strengthening the sandbox environment will help maintain system integrity and protect underlying infrastructure.

\textbf{Discussion System Enhancement:} While the platform supports CRUD operations for discussions and comments, additional features such as tagging, advanced search functionality, and content moderation can be introduced. In particular, enhancing the search mechanism with keyword indexing, filtering options (e.g., by topic, popularity, or recency), and relevance ranking would significantly improve information retrieval. These improvements would enable users to more easily locate related threads, explore previously discussed topics, and engage with content that aligns with their interests, thereby fostering a more organized and accessible knowledge-sharing environment.

\textbf{Multi-language Support Expansion:} The current system supports a limited set of programming languages, including C\#, Java and Python. Expanding support to a wider range of languages—such as Go, Rust, and others—would significantly enhance the platform’s accessibility and appeal to a broader audience. This improvement would allow users with different programming backgrounds to practice more effectively and better reflect real-world scenarios where multiple languages are commonly used. Additionally, supporting diverse languages would increase the system’s flexibility and position it as a more comprehensive coding practice platform.

\textbf{CI/CD Integration:} 

To further improve development efficiency, system reliability, and deployment consistency, the integration of a Continuous Integration and Continuous Deployment (CI/CD) pipeline is a crucial next step for the platform. Currently, the development workflow relies on manual processes for testing and deployment, which can introduce delays and increase the risk of human error. Implementing a CI/CD pipeline would automate these processes, enabling faster and more reliable software delivery.


In summary, these proposed improvements aim to enhance the platform across multiple dimensions, including user engagement, system performance, security, and intelligent features. Collectively, they provide a clear direction for evolving the system into a more robust, scalable, and user-centric coding practice platform.
\newpage